%%% Praxis: Formally Verified Persistent Memory for Agentic AI
%%% ACM acmart (sigplan) — draft for Coblenz demo, July 2026
\documentclass[sigplan,screen,nonacm]{acmart}

\usepackage{listings}
\usepackage{xcolor}
\usepackage{booktabs}
\usepackage{subcaption}

% F* convenience macros
\newcommand{\fstar}{F\ensuremath{^\star}}
\newcommand{\pointsto}{\ensuremath{\mapsto}}
\newcommand{\sepconj}{\ensuremath{\mathbin{*\!*}}}

% Listing style for F*
\lstdefinelanguage{fstar}{
  morekeywords={module,open,type,let,fn,val,match,with,if,else,not,
    requires,ensures,returns,pure,Pure,erased,nat,bool,string,true,false,
    Pulse,Lib,Pervasives,FStar,List,Tot,String},
  sensitive=true,
  morecomment=[l]{//},
  morecomment=[s]{(*}{*)},
  morestring=[b]",
  literate={|->}{{$\mapsto$}}2
           {**}{{$\mathbin{*\!*}$}}2
           {/\\}{{$\wedge$}}2
           {<=}{{$\leq$}}2
           {=>}{{$\Rightarrow$}}2,
}

\lstset{
  basicstyle=\small\ttfamily,
  keywordstyle=\bfseries\color{blue!70!black},
  commentstyle=\itshape\color{gray},
  stringstyle=\color{red!70!black},
  columns=flexible,
  keepspaces=true,
  frame=single,
  framerule=0.4pt,
  xleftmargin=1em,
  xrightmargin=0.5em,
  aboveskip=0.8em,
  belowskip=0.5em,
  breaklines=true,
  showstringspaces=false,
  tabsize=2,
  captionpos=b,
}

%%% ---------- metadata ----------

\title{Praxis: Formally Verified Persistent Memory for Agentic AI}
\subtitle{Preventing Hallucination Persistence Loops with Separation Logic on the Red Hat AI Stack}

\author{Gerald Trotman}
\affiliation{%
  \institution{Red Hat}
  \country{USA}
}

\begin{abstract}
AI agents with persistent memory face a class of errors no existing
framework addresses with formal methods: an agent hallucinates, saves the
hallucination to persistent storage, and treats its own fiction as ground
truth in all future sessions---a self-reinforcing \emph{hallucination
persistence loop}. We present Praxis ($\pi\rho\hat{\alpha}\xi\iota\varsigma$---verified
action), a runtime verification system that imposes a proof obligation on
every persist operation across five agent execution surfaces: model
inference, code generation, tool invocation, skill persistence, and
multi-agent communication. Praxis uses \fstar{}/Pulse separation logic with
Z3 SMT proof discharge to verify 10 properties spanning content soundness
and memory substrate safety. We deploy Praxis as a verification sidecar
alongside the Hermes agentic framework on Red Hat OpenShift AI with
vLLM/KServe model serving and NVIDIA OpenShell sandbox enforcement. Our
evaluation shows 22 verified \fstar{} specifications, 140 Python runtime
tests, and zero-admit formal proofs covering OWASP ASI06 memory poisoning
vectors. To our knowledge, Praxis is the first system to apply separation
logic to AI agent persistent memory.
\end{abstract}

\keywords{formal verification, separation logic, agentic AI, persistent memory, hallucination detection, \fstar{}, Pulse}

\begin{document}

\maketitle

%%% ================================================================
\section{Introduction}
\label{sec:intro}

Agentic AI frameworks---systems where an LLM autonomously invokes tools,
generates code, and persists state across sessions---are moving into
production. Hermes~\cite{hermes2026} (Nous Research), OpenClaw, and Kagent
(Red Hat) each implement persistent memory architectures that survive pod
restarts, enabling agents to learn from experience. This persistence is also
the attack surface.

\paragraph{The hallucination persistence loop.}
An agent observes ``GPU utilization at 87\%,'' but the model confabulates
``the cluster is overloaded and will crash.'' If this conclusion enters
persistent memory, every future session begins with ``the cluster is
overloaded'' as established fact. The agent compounds the error: it
references the fabricated fact, derives further conclusions, and stores those
too. Each session reinforces the fiction.

No existing mitigation provides formal guarantees against this:

\begin{itemize}
  \item Sandbox isolation (OpenShell~\cite{nvidia2026openshell}, gVisor)
    constrains runtime behavior but does not inspect content.
  \item Content filtering (heuristic similarity, embedding distance) catches
    some fabrications but cannot prove soundness.
  \item Type systems for agent control
    (LBAC/TypeGuard~\cite{zhou2026lbac}, TACIT~\cite{odersky2026tacit})
    verify code safety but do not address persistent memory.
  \item Verification-aware code generation
    (Dafny-as-IL~\cite{li2025dafny}) verifies generated code but not stored
    reasoning.
\end{itemize}

\textbf{Praxis} closes this gap by treating every memory write as a
verification obligation. The agent must prove---using a formally verified
pipeline---that what it persists is derivable from real observations through
valid inference. Not ahead-of-time program verification, but
\emph{proof-per-action at runtime} across five execution surfaces.

\subsection{Five Verification Surfaces}

\begin{table}[t]
\caption{Praxis verification surfaces and their correspondence to Hermes agent features.}
\label{tab:surfaces}
\begin{tabular}{@{}lll@{}}
\toprule
\textbf{Surface} & \textbf{What Praxis Proves} & \textbf{Hermes Feature} \\
\midrule
1. Model inference & Conclusion follows from observations & vLLM/KServe $\to$ reasoning \\
2. Code generation & Generated skill satisfies spec & Auto-generated skills \\
3. Tool invocation & Tool output is trustworthy & kubectl, web scrape, API \\
4. Skill persistence & Stored skill still valid & MEMORY.md, SQLite, skills \\
5. Multi-agent & Agent B can verify Agent A's proof & Multi-platform sessions \\
\bottomrule
\end{tabular}
\end{table}

\subsection{Contributions}

\begin{enumerate}
  \item \textbf{Formalization.} 10 verification properties (P1--P10)
    spanning content soundness and memory substrate safety, specified in
    \fstar{}/Pulse with Z3 proof discharge. The first application of
    separation logic to AI agent persistent memory.

  \item \textbf{Verification strength spectrum.} An auto-active verification
    design~\cite{leino2010dafny} that balances expressiveness and control:
    five inference rule categories from Identity (full verification) through
    Derivation (weakest). The agent chooses the rule; the system checks the
    obligation.

  \item \textbf{Incorrectness lemmas.} Following
    Gardner's~\cite{gardner2026incorrectness} incorrectness logic, we prove
    not only that verified writes satisfy properties (soundness of
    acceptance) but that fabricated inputs are necessarily rejected
    (soundness of rejection).

  \item \textbf{Deployable architecture.} Praxis deploys as a verification
    sidecar on Red Hat OpenShift AI alongside Hermes with vLLM/KServe model
    serving and NVIDIA OpenShell~\cite{nvidia2026openshell} runtime sandbox.
    KaRaMeL extraction~\cite{protzenko2017verified} compiles verified
    \fstar{}/Pulse to native C for production latency targets.
\end{enumerate}


%%% ================================================================
\section{Background and Motivation}
\label{sec:background}

\subsection{Agent Persistent Memory Architectures}

\textbf{Hermes Agent}~\cite{hermes2026} (Nous Research) implements a closed
learning loop on Red Hat OpenShift AI: it generates reusable skills from
multi-step tasks, maintains bounded persistent memory (MEMORY.md at 2{,}200
chars, USER.md at 1{,}375 chars), and searches conversation history via
SQLite FTS5. Skills self-improve based on feedback and persist across pod
restarts via PersistentVolumeClaim. Deployed on UBI~9 with vLLM model
serving via KServe InferenceService.

\paragraph{The persistence surface.}
Hermes writes to five targets: curated memory (MEMORY.md), user model
(USER.md), conversation database (SQLite), generated skills
(\texttt{/opt/data/skills/}), and cron schedules. Every write survives pod
restarts. A fabrication entering any target compounds indefinitely.

\subsection{Related Verification Approaches}

\textbf{Dafny as Verification-Aware Intermediate Language}~\cite{li2025dafny}.
LLM generates Dafny code $\to$ Dafny verifier checks $\to$ compiled to
target language. Addresses Surface~2 (code generation) but not persistent
memory. Praxis adopts the ``verification-aware IL'' pattern: \fstar{} is our
intermediate verification language, but we verify \emph{inference
soundness}, not just code correctness.

\textbf{Language-Based Agent Control / TypeGuard}~\cite{zhou2026lbac}.
Policies encoded as Haskell types; agent code must type-check before
execution. Abstract data types enforce data provenance and information flow
control. Dual LLM architecture (privileged + quarantined) generalizes
capability separation. Addresses Surfaces~1 and~3. Praxis adopts the
``policies-as-types'' pattern: our algebraic \texttt{inference\_rule} type
forces evidence per claim. But Praxis uses refinement types with proof
discharge (stronger than Haskell's type system) and addresses persistent
memory (which LBAC does not).

\textbf{Tracked Capabilities for Safer Agents / TACIT}~\cite{odersky2026tacit}.
Scala~3 capture checking tracks capabilities through the type system.
\texttt{Classified[T]} wrapper enforces local purity---functions processing
classified data cannot leak it. MCP server implementation. Addresses
Surfaces~3 and~5. Praxis adopts the ``tracked capabilities'' pattern: proof
certificates are capabilities that must be presented on read. But Praxis
uses separation logic for memory ownership (P5, P6)---a dimension no other
agent verification system addresses.

\subsection{What Praxis Adds Beyond All Three}

\begin{table}[t]
\caption{Capability comparison with concurrent verification approaches for agentic AI.}
\label{tab:comparison}
\begin{tabular}{@{}lcccc@{}}
\toprule
\textbf{Capability} & \textbf{Dafny-IL} & \textbf{LBAC} & \textbf{TACIT} & \textbf{Praxis} \\
\midrule
Code verification        & \checkmark & ---        & ---        & \checkmark \\
Inference soundness      & ---        & \checkmark & ---        & \checkmark \\
Tool trust               & ---        & \checkmark & \checkmark & \checkmark \\
Capability tracking      & ---        & ---        & \checkmark & \checkmark \\
Multi-agent safety       & ---        & ---        & \checkmark & \checkmark \\
\textbf{Persistent memory}       & ---        & ---        & ---        & \textbf{\checkmark} \\
\textbf{Separation logic}        & ---        & ---        & ---        & \textbf{\checkmark} \\
\textbf{Memory poisoning (ASI06)} & ---       & ---        & ---        & \textbf{\checkmark} \\
\textbf{Red Hat stack}            & ---        & ---        & ---        & \textbf{\checkmark} \\
\bottomrule
\end{tabular}
\end{table}


%%% ================================================================
\section{Formalization}
\label{sec:formalization}

\subsection{Core Types (\texttt{PraxisTypes.fst})}

Algebraic type with evidence per inference rule---the auto-active pattern:

\begin{lstlisting}[language=fstar,caption={Inference rule algebraic type with per-variant evidence.},label={lst:rule}]
type inference_rule =
  | RuleIdentity
  | RuleExtraction  : ev:extraction_evidence  -> inference_rule
  | RuleAggregation : ev:aggregation_evidence -> inference_rule
  | RuleToolResult  : ev:tool_evidence        -> inference_rule
  | RuleDerivation  : ev:derivation_evidence  -> inference_rule
\end{lstlisting}

Graduated trust for observation provenance:

\begin{lstlisting}[language=fstar,caption={Trust levels for observation sources.},label={lst:trust}]
type trust_level =
  | DirectObservation | ToolOutput | ExternalInput
  | AgentGenerated | Unverified
\end{lstlisting}

\subsection{Verification Strength Spectrum (\texttt{PraxisPredicates.fst})}

\begin{table}[t]
\caption{Verification strength spectrum: the agent chooses the rule; the system checks the obligation.}
\label{tab:spectrum}
\begin{tabular}{@{}llll@{}}
\toprule
\textbf{Rule} & \textbf{Obligation} & \textbf{Strength} & \textbf{Expressiveness} \\
\midrule
Identity    & conclusion = single premise   & Full     & Observation relay only \\
Extraction  & conclusion $\subseteq$ source & Strong   & Substring extraction \\
Aggregation & $|$conclusion$|$ $\leq$ $\Sigma|$premises$|$ & Moderate & Multi-source synthesis \\
ToolResult  & tool is trusted               & Delegated & Full tool output \\
Derivation  & rule\_id non-empty, conf $\in$ [0,100] & Weakest & Domain reasoning \\
\bottomrule
\end{tabular}
\end{table}

\subsection{Content Properties (P1--P4)}

\begin{description}
  \item[P1 Inference Soundness] (\texttt{AgentReasoning}):
    \texttt{chain\_well\_formed} checks structural ordering AND semantic rule
    obligations.
  \item[P2 Consistency:] \texttt{no\_contradiction} against existing memory
    state.
  \item[P3 Poisoning Resistance] (\texttt{PoisonDetection}):
    Case-insensitive pattern matching against OWASP
    ASI06~\cite{owasp2026asi06} vectors---instruction overrides, role
    escalations, exfiltration markers.
  \item[P4 Completeness] (\texttt{CompletenessCheck}): Domain-spec critical
    observations covered and required keys present.
\end{description}

\subsection{Substrate Properties (P5--P8)}

Pulse separation logic for memory ownership:

\begin{lstlisting}[language=fstar,caption={Verified write with Pulse separation logic postconditions.},label={lst:vwrite}]
fn verified_write (agent: agent_id) (region: memory_region) ...
  requires region |-> v
  returns r: write_result
  ensures (match r with
           | WriteOk -> region |-> conclusion.me_content
                        ** pure (...)
           | _ -> region |-> v)
\end{lstlisting}

\begin{description}
  \item[P5 Ownership:] \texttt{region \pointsto{} v}---agent holds the
    points-to assertion.
  \item[P6 Frame Rule:] Writing to MEMORY.md does not modify USER.md
    (separation logic frame).
  \item[P7 Bounds:] Content length $\leq$ tier bound (2{,}200 chars for
    Hermes Tier~1).
  \item[P8 Concurrency:] Multi-session access to shared SQLite is race-free.
\end{description}

\subsection{Tool and Skill Properties (P9--P10)}

\begin{description}
  \item[P9 Tool Scope:] Tool call arguments within declared API surface.
  \item[P10 Side-Effect Containment:] Observed side effects $\subseteq$
    declared side effects.
\end{description}

\subsection{Skill Verification (Surface 2)}

Generated skills must satisfy a specification before persisting:

\begin{lstlisting}[language=fstar,caption={Skill safety predicate composing interface, scope, and content checks.},label={lst:skill}]
val skill_safe_to_persist :
  spec:skill_spec -> code:skill_code -> pp:poison_patterns ->
  Pure bool
    (ensures fun b -> b ==>
      skill_interface_valid spec code /\
      skill_tools_within_scope spec code /\
      content_safe code.sc_body pp)
\end{lstlisting}

\subsection{Temporal Validity (Surface 4)}

Proof certificates carry validity windows for re-verification:

\begin{lstlisting}[language=fstar,caption={Temporal validity: certificates expire and must be re-verified.},label={lst:temporal}]
val certificate_valid_at :
  cert:temporal_certificate -> time:nat ->
  Pure bool
    (ensures fun b -> b ==>
      cert.tc_valid_from <= time /\
      time <= cert.tc_valid_until)
\end{lstlisting}

\subsection{Multi-Agent Proof Transport (Surface 5)}

Session types~\cite{gay2026session} for cross-agent proof witness transport:

\begin{lstlisting}[language=fstar,caption={Cross-agent verification: Agent B verifies Agent A's proof witness.},label={lst:witness}]
val cross_agent_verify :
  witness:proof_witness -> entry:memory_entry ->
  Pure bool
    (ensures fun b -> b ==>
      witness.pw_content_hash = hash entry.me_content /\
      witness.pw_source_trusted)
\end{lstlisting}

\subsection{Normalization Tests (\texttt{PraxisNormTests.fst})}

\texttt{assert\_norm} evaluates verified functions on concrete inputs at
type-checking time---pure computation in the \fstar{} normalizer, no Z3
needed. Bridges universal proofs and existential tests by confirming the
specification computes correctly on the same inputs the Python test suite
exercises.

\subsection{Incorrectness Lemmas (\texttt{PraxisLemmas.fst})}

Following Gardner~\cite{gardner2026incorrectness} (OPLSS 2026), we prove
rejection is \emph{necessary} under attack conditions:

\begin{itemize}
  \item L1: Identity fabrication $\to$ always rejected
  \item L2: Multi-premise identity $\to$ always rejected
  \item L3: Untrusted tool output $\to$ always rejected
  \item L4: Empty inference chain $\to$ always rejected
  \item L5: Evidence-free derivation $\to$ always rejected
  \item L6: Over-confident derivation $\to$ always rejected
  \item L7: Contradiction with existing memory $\to$ always rejected
\end{itemize}


%%% ================================================================
\section{Implementation}
\label{sec:implementation}

\subsection{Praxis Verification Gateway}

\paragraph{Python runtime architecture.}
The Python runtime mirrors the \fstar{} specification with structural
fidelity. \texttt{gateway\_client.py} (533 lines) implements
\texttt{PraxisGatewayClient}, which executes the same verification pipeline
as \texttt{VerifiedWrite.fst}: P1+P2 (inference soundness and consistency)
followed by P3 (poison detection) followed by P4 (completeness) followed by
P5+P7 (ownership and bounds). Each \fstar{} function---\texttt{chain\_well\_formed},
\texttt{no\_contradiction}, \texttt{content\_safe},
\texttt{completeness\_check}---has a Python mirror that accepts the same
arguments, applies the same case analysis, and returns the same verdict. The
algebraic \texttt{inference\_rule} type maps to a Python dataclass hierarchy
with identical constructors: \texttt{RuleIdentity},
\texttt{RuleExtraction(evidence)}, \texttt{RuleAggregation(evidence)},
\texttt{RuleToolResult(evidence)}, and \texttt{RuleDerivation(evidence)}.

\paragraph{gRPC service definition.}
\texttt{praxis.proto} defines six RPCs that cover all five verification
surfaces: \texttt{VerifyWrite} (Surface~1, general memory persistence),
\texttt{VerifyInference} (Surface~1, inference chain verification),
\texttt{VerifySkill} (Surface~2, skill persistence),
\texttt{VerifyToolInvocation} (Surface~3, tool output trust),
\texttt{VerifyTemporal} (Surface~4, re-verification with validity windows),
and \texttt{VerifyWitness} (Surface~5, cross-agent proof transport). The
Protobuf \texttt{oneof} construct in the \texttt{InferenceRule} message
mirrors the \fstar{} algebraic \texttt{inference\_rule} type: each variant
carries its own evidence payload, ensuring that the wire format enforces the
same structural obligations as the formal specification.

\paragraph{Proof certificate lifecycle.}
On a successful \texttt{WriteOk} verdict, the gateway constructs a proof
certificate containing: a SHA-256 content hash of the persisted content, the
inference rule that justified the write, the list of properties satisfied (a
subset of P1--P10), and an OpenTelemetry trace ID linking the certificate to
the observability pipeline. For Surface~4, a \texttt{TemporalCertificate}
wraps the base certificate with a validity window (\texttt{valid\_from},
\texttt{valid\_until}), enabling periodic re-verification. For Surface~5, a
\texttt{ProofWitness} packages the certificate with the originating agent's
identity and a trust attestation.

\paragraph{Sidecar deployment.}
\texttt{hermes-sidecar-patch.yaml} patches the Hermes agent pod to include
the Praxis verification sidecar as an additional container. The sidecar
listens on port 50052 and forwards verified requests to the gateway on port
50051. Setting \texttt{HERMES\_VERIFY\_WRITES=true} enables write
interception. The entire sidecar runs within the NVIDIA
OpenShell~\cite{nvidia2026openshell} sandbox.

\subsection{Deployment on Red Hat OpenShift AI}

\begin{lstlisting}[language={},caption={Deployment architecture: Hermes agent with Praxis sidecar on OpenShift AI.},label={lst:arch}]
+- OpenShift AI ----------------------------------+
|                                                  |
|  vLLM/KServe InferenceService (GPU serving)     |
|         |                                        |
|         v  OpenAI-compatible API                 |
|  +- NVIDIA OpenShell Sandbox -----------------+ |
|  |                                             | |
|  |  +- Hermes Agent (UBI 9) ---------------+  | |
|  |  |  MEMORY.md, USER.md, SQLite, skills  |  | |
|  |  +----------------+---------------------+  | |
|  |                   | gRPC MemoryIntent       | |
|  |  +----------------v---------------------+  | |
|  |  |  Praxis Sidecar                      |  | |
|  |  |  F*/Z3 -> certificate or rejection   |  | |
|  |  |  OTel span -> MLflow audit trail     |  | |
|  |  +--------------------------------------+  | |
|  |  Policy Engine (seccomp, Landlock, eBPF)    | |
|  +---------------------------------------------+ |
|                                                  |
|  PostgreSQL (CNPG) | OTel | MLflow | Tempo      |
+--------------------------------------------------+
\end{lstlisting}

\subsection{Complementary Enforcement}

\begin{table}[t]
\caption{Division of enforcement between OpenShell (runtime sandbox) and Praxis (content verification).}
\label{tab:enforcement}
\begin{tabular}{@{}lcc@{}}
\toprule
\textbf{Concern} & \textbf{OpenShell} & \textbf{Praxis} \\
\midrule
Agent can't read another's files & Policy Engine & P5 ownership \\
Memory write within bounds       & Policy Engine & P7 bounds \\
Agent can't exfiltrate           & Policy Engine & --- \\
Stored reasoning is sound        & ---           & P1 soundness \\
Memory not poisoned              & ---           & P3 (ASI06) \\
Sensitive context stays local    & Privacy Router & --- \\
Concurrent memory access         & ---           & P8 concurrency \\
\bottomrule
\end{tabular}
\end{table}


%%% ================================================================
\section{Evaluation}
\label{sec:evaluation}

\subsection{Verification Results}

\paragraph{Methodology.}
\fstar{} specifications are verified using \texttt{fstar.exe} in dependency
order: \texttt{PraxisTypes.fst} (no dependencies), then
\texttt{PraxisPredicates.fst}, then property modules, then composite
modules (\texttt{VerifiedWrite.fst}, \texttt{VerifiedSkillWrite.fst}), and
finally test and lemma modules. Python runtime tests use \texttt{pytest}
with markers for \fstar{} mirror tests, OWASP vectors, and integration
tests.

\begin{table}[t]
\caption{Verification evidence summary.}
\label{tab:results}
\begin{tabular}{@{}lrl@{}}
\toprule
\textbf{Artifact} & \textbf{Count} & \textbf{Status} \\
\midrule
\fstar{} specifications    & 22 & All verified (zero admits) \\
Python runtime tests        & 140 & All pass \\
OWASP ASI06 vectors         & 23 + 8 benign & Full coverage \\
Normalization tests         & 20 assert\_norm & Type-checker evaluated \\
Incorrectness lemmas        & 7 & All verified \\
Proto messages              & 30+ & Request/response/enum \\
Hallucination loop          & 1 canonical & Rejected at P1 \\
\bottomrule
\end{tabular}
\end{table}

\subsection{Verification Latency}

\fstar{} verification is a build-time cost, not a runtime cost. Each
\fstar{} module verifies in 1--3 seconds; the full 22-module suite
completes in under 40 seconds. The runtime Python gateway performs pure
function evaluation---pattern matching, string comparison, length checks,
hash computation---on bounded inputs. Hermes Tier~1 content (MEMORY.md) is
capped at 2{,}200 characters; Tier~2 (USER.md) at 1{,}375 characters. On
these bounded inputs, the Python verification pipeline completes in
microseconds. KaRaMeL extraction to C~\cite{protzenko2017verified} would
reduce even this cost to sub-microsecond levels.

\subsection{OWASP ASI06 Attack Coverage}

The OWASP ASI06~\cite{owasp2026asi06} (Memory Poisoning) attack surface is
tested with 31 vectors organized into four categories.
\emph{Instruction overrides} (8 vectors): attempts to inject system-level
directives into persistent memory. \emph{Role escalation} (6 vectors):
attempts to elevate trust level. \emph{Exfiltration markers} (6 vectors):
attempts to embed data exfiltration payloads. \emph{Compound attacks} (3
vectors): multi-technique payloads. All 23 attack vectors are rejected at P3
by \texttt{PoisonDetection.fst}. Additionally, 8 benign vectors are
confirmed as accepted, establishing that the detector is false-negative-free
without being overly restrictive.

\subsection{Normalization Test Bridge}

The \texttt{assert\_norm} mechanism in \fstar{} evaluates a verified
function on a concrete input during type-checking---the \fstar{} normalizer
reduces the expression to a boolean value and the type-checker confirms it
equals \texttt{true}. The normalization test bridge exploits this to connect
two distinct verification regimes: universal proofs (Z3 discharges a
property for ALL valid inputs) and existential tests (pytest confirms the
property for THIS specific input on the Python mirror). If the \fstar{} spec
and the Python mirror ever diverge, the bridge catches it.

\subsection{Incorrectness Lemma Coverage}

Following Gardner's incorrectness logic~\cite{gardner2026incorrectness},
\texttt{PraxisLemmas.fst} contains 7 lemmas (L1--L7) that prove rejection
is \emph{necessary} under specific attack conditions, not merely that
acceptance is correct.

\begin{table}[t]
\caption{Incorrectness lemmas: each proves that a specific attack class is necessarily rejected.}
\label{tab:lemmas}
\begin{tabular}{@{}lll@{}}
\toprule
\textbf{Lemma} & \textbf{Attack condition} & \textbf{Rejection} \\
\midrule
L1 & Identity with fabricated conclusion & P1 rejects \\
L2 & Identity with multiple premises & P1 rejects \\
L3 & ToolResult with untrusted source & P1 rejects \\
L4 & Empty inference chain & P1 rejects \\
L5 & Derivation with empty rule id & P1 rejects \\
L6 & Derivation with confidence $>$ 100 & P1 rejects \\
L7 & Contradiction with existing memory & P2 rejects \\
\bottomrule
\end{tabular}
\end{table}

L5 is particularly significant: even the weakest inference rule (Derivation)
rejects evidence-free claims. An agent cannot bypass the verification gate
by choosing the least restrictive rule---every rule demands \emph{some}
evidence.


%%% ================================================================
\section{Discussion}
\label{sec:discussion}

\subsection{Limitations}

\paragraph{(a) The Derivation rule is intentionally weak.}
The weakest inference rule, \texttt{RuleDerivation}, requires only a
non-empty domain rule identifier and a confidence score in [0, 100]. It does
not verify the validity of the domain rule itself. This is a deliberate
design choice: domain-specific abductive reasoning cannot be verified by a
general-purpose formal system without a domain ontology. Strict deployments
can extend Praxis with a \emph{domain rule catalog}. Crucially, incorrectness
lemma L5 proves that even this weakest rule rejects evidence-free claims.

\paragraph{(b) Specification correctness.}
Formal verification proves that code meets its specification. It does not
prove that the specification captures the intended behavior. If the poison
detection patterns miss a novel attack vector, the verifier will accept it
as benign. This is a fundamental limitation of all formal
methods~\cite{rand2026verified}. The normalization test bridge (Section~\ref{sec:evaluation})
partially mitigates this.

\paragraph{(c) KaRaMeL extraction.}
The current implementation uses a Python mirror of the \fstar{}
specification. The production target is KaRaMeL
extraction~\cite{protzenko2017verified}: the \fstar{}/Pulse specifications
compile to C via the KaRaMeL compiler, producing a verified native library.
This eliminates the mirror-consistency question entirely.

\subsection{Comparison to the Dual LLM Pattern}

LBAC/TypeGuard~\cite{zhou2026lbac} proposes a binary separation: a
privileged LLM with tool access that never processes untrusted data, and a
quarantined LLM that processes untrusted data but has no tool access.
TACIT~\cite{odersky2026tacit} refines this with Scala~3 capture checking:
\texttt{Classified[T]} wraps values so that functions processing classified
data cannot leak them.

Praxis generalizes beyond the binary privileged/quarantined distinction with
a five-level trust spectrum. Rather than classifying the \emph{agent} as
privileged or quarantined, Praxis classifies each \emph{inference step} on a
continuous strength scale from Identity (strongest) through Derivation
(weakest). This graduated model accommodates the reality that agent
reasoning involves a mix of trust levels within a single session.

\subsection{Separation Logic for Agent Memory}

Praxis makes a novel contribution by applying separation logic to AI agent
persistent memory. Pulse---the separation logic layer of \fstar{}---provides
the right formalism for this domain. The points-to assertion
\texttt{region \pointsto{} v} models exclusive ownership of a memory region
by an agent, directly encoding the Hermes memory architecture. The
separating conjunction \sepconj{} guarantees write isolation:
\texttt{region\_a \pointsto{} v\_a \sepconj{} region\_b \pointsto{} v\_b}
proves that a write to \texttt{region\_a} cannot affect
\texttt{region\_b}---the formal statement of property P6 (frame rule).

The \texttt{stt} (stateful) effect in Pulse sequences stateful operations
with explicit pre- and post-conditions. Erased ghost state (\texttt{erased}
types) tracks preconditions and invariants without runtime cost: the
ownership proof exists at verification time but compiles away to zero
overhead in the extracted code.


%%% ================================================================
\section{Conclusion}
\label{sec:conclusion}

Praxis demonstrates that formal verification of AI agent persistent memory
is both tractable and deployable. By treating every memory write as a proof
obligation---verified by \fstar{}/Pulse separation logic and Z3 SMT
discharge---we prevent hallucination persistence loops at the formal methods
level, not the heuristic level.

We claim four contributions. First, Praxis is the first application of
separation logic to AI agent persistent memory: the points-to assertion
models agent ownership, the separating conjunction guarantees write
isolation, and the frame rule proves that skill persistence does not corrupt
curated memory. Second, the verification strength spectrum provides five
inference rule categories that give deployers a graduated trust model rather
than binary privileged/quarantined separation. Third, incorrectness lemmas
prove that rejection is necessary under attack conditions. Fourth, the
architecture is deployable today on Red Hat OpenShift AI with Hermes,
vLLM/KServe, and OpenShell.

Praxis complements rather than competes with the three concurrent
verification approaches. Dafny-as-IL~\cite{li2025dafny} verifies generated
code. LBAC/TypeGuard~\cite{zhou2026lbac} encodes agent policies as types.
TACIT~\cite{odersky2026tacit} tracks capabilities. Praxis adds persistent
memory verification---the dimension that none of the three address.

Future work proceeds along four axes. First, KaRaMeL
extraction~\cite{protzenko2017verified} will compile the \fstar{}/Pulse
specifications to C, eliminating the Python mirror. Second, an extended
domain rule catalog for the Derivation inference rule will reduce the
weakest rule to a lookup against curated reasoning patterns. Third,
multi-agent proofs with the Kagent (Red Hat) fleet will extend Surface~5
using session types~\cite{gay2026session}. Fourth, we are investigating
Praxis as a compliance evidence generator for EU AI Act Article~9 (risk
management), where proof certificates and OTel audit trails provide
machine-readable evidence of verification.


%%% ================================================================
\bibliographystyle{ACM-Reference-Format}
\bibliography{references}

\end{document}
